\chapter{アルゴリズムの基礎}
\label{c-algo-basics}

本章では、式~\eqref{eq-mk-discr}の主問題またはその双対~\eqref{eq-dual}として記述される optimal transport の離散定式化を解くために使用できる、組合せ最適化および線形計画法における最も一般的なアルゴリズム的ツールを説明する。

これらのアルゴリズムの起源は第二次世界大戦にまで遡ることができ、\citeauthor{tolstoi1930methods}の先駆的研究\citeyearpar{tolstoi1930methods}による戦前の業績か、あるいは戦時中に~\citet{Hitchcock41}と\citet{Kantorovich42}が利用可能な資源を消費地点に最適に配分するという一般的な問題を定式化した時期に始まる。これらの定式化、および後の~\citet{koopmans1949optimum}の貢献はいずれも、その問題を解くための\emph{証明可能に}正しいアルゴリズムを提供するには至らなかった（巡回違反法は~\citet{tolstoi1939metody}によって既にヒューリスティックとして提案されていた）。線形計画法の分野が完全に開花し、シンプレックス法が提案されるまで、これらの問題を厳密に解くことはできなかった。

線形計画法の目的は、目的関数が線形であり、制約が関心のある変数に関する線形（不）等式であるような最適化問題を解くことである。Optimal transport 問題はその記述に適合し、したがってこのより広いクラスの問題の特殊な場合である。しかしながら、optimal transport はすべての線形計画の中で真に特別であると主張できる。第一に、\citeauthor{dantzig1991}の線形計画を解くための初期の動機は、輸送問題を解くことと大いに関連していた~\citep[p. 210]{dantzig49econometrica}。第二に、特殊な場合に過ぎないにもかかわらず、optimal transport 問題は最適化の脚光を浴び続けた。それは、optimal transport 問題が~\citet{ford1962flows}の結果により、最小費用ネットワークフローとして知られる線形計画の重要なクラスと関連し、実際に等価であることがほどなく理解されたためである~\citep[p. 213, Lem. 9.3]{korte2012combinatorial}。このように、OT 問題は数理計画法の誕生以来~\citep{Dantzig51}特別な注目を受けてきており、今でも新しい聴衆に最適化を紹介するために広く使用されている~\citep[\S1, p. 4]{nocedal}。

\section{Kantorovich の線形計画}\label{s-kantorovich}
式~\eqref{eq-mk-discr}において主 OT 問題を既に導入した：
\eql{\label{eq-mk-discr-algo}
	\MKD_{\C}(\a,\b) =
	\umin{\P \in \CouplingsD(\a,\b)}
		\sum_{i\in \range{n}, j \in \range{m}} \C_{i,j} \P_{i,j}.
}
線形計画法の文献との関連を明確にするために、上記の方程式を\emph{標準}形の線形計画、すなわち線形目的関数、行列と定数ベクトルで定義される等式制約、および変数の非負制約を持つ線形計画として書き直すことができる。$\Identity_n$ をサイズ $n$ の単位行列、$\otimes$ を Kronecker 積とする。${(n+m) \times nm}$ 行列
$$
\mathbf{A}= \begin{bmatrix}
	\transp{\ones_{n}} \otimes \Identity_m \\
	\Identity_n \otimes \transp{\ones_{m}}
\end{bmatrix}\in\RR^{(n+m)\times nm}
$$
を用いて、任意の $\P$ が $\CouplingsD(\a,\b)$ に属するために満たすべき行和制約と列和制約を符号化できる。そのために、行列 $\P\in\RR^{n\times m}$ をベクトル $\p\in\RR^{nm}$ として書き直す。ここで $\p$ の $i+n(j-1)$ 番目の要素は $\P_{ij}$ に等しい（$\P$ は列方向に列挙される）。すると以下の同値性が得られる：
$$\P\in \RR^{n\times m}\in \CouplingsD(\a,\b) \Leftrightarrow \p\in\RR^{nm}_+, \mathbf{A}\p = \bigl[\begin{smallmatrix}\a\\ \b \end{smallmatrix} \bigr].$$
したがって、元の optimal transport 問題を次のように書くことができる：
\eql{\label{eq-mk-discr-primal}
	\MKD_{\C}(\a,\b) =
	\umin{\substack{\p \in \RR^{nm}_+\\
		\mathbf{A}\p = \bigl[\begin{smallmatrix}\a\\ \b \end{smallmatrix} \bigr]}} \transp{\cc}\p,
}
ここで $nm$ 次元ベクトル $\cc$ はコスト行列 $\C$ に含まれる列を積み重ねたものに等しい。

\begin{rem}\label{rem-transportation-polytope} 上記の $n+m$ 個の制約のうち1つは冗長であること、言い換えれば、行列 $A$ の行ベクトルは線形独立でないことに注意せよ。実際、最初の $n$ 行とその後の $m$ 行をそれぞれ合計すると、同じベクトルが得られる（すなわち $A \bigl[\begin{smallmatrix} \ones_n\\\zeros_m \end{smallmatrix}\bigr] = A \bigl[\begin{smallmatrix} \zeros_n\\\ones_m\end{smallmatrix}\bigr]=\transp{\ones_{nm}}$）。$A$ から1行を除き、$\bigl[\begin{smallmatrix}\a\\ \b \end{smallmatrix}\bigr]$ の対応する成分を除いても、適切に定義された線形系が得られることを示すことができる。簡単のため、また $\a$ と $\b$ を非対称に扱うことを避けるため、以下では冗長な定式化を保持する。ただし、退化がいくつかの計算で現れることに留意されたい。\end{rem}

式~\eqref{eq-mk-discr-primal}に対応する双対問題は、線形計画法における双対性~\citep[p. 143]{bertsimas1997introduction}に従い、次のように定義される：
\eql{\label{eq-mk-discr-dual}
	\MKD_{\C}(\a,\b) =
	\umax{\substack{\hD \in \RR^{n+m}\\
		\transp{A}\hD \leq \cc}}  \transp{\bigl[\begin{smallmatrix}\a\\ \b \end{smallmatrix} \bigr]}\hD.
}

この計画は式~\eqref{prop-duality-discr}で提示されたものと正確に等価であることに注意せよ。


\begin{rem}\label{rem-duality} 上記の双対性の結果について、注意~\ref{rem-kantorovich-dual}で展開された議論の直接的な定式化と見なすことができる簡単な導出を与える。強双対性、すなわち主問題~\eqref{eq-mk-discr-primal}と双対問題~\eqref{eq-mk-discr-dual}の最適値が実際に一致するという事実は、より長い証明を必要とする~\citep[\S4.10]{bertsimas1997introduction}。記法を簡略化するために、$\q=\bigl[\begin{smallmatrix}\a\\ \b \end{smallmatrix} \bigr]$ と書く。
ここで optimal transport 問題の緩和された主問題を考える。制約 $\mathbf{A}\p=\q$ はもはや必ずしも課されず、代わりに任意のコストベクトル $\hD\in\RR^{n+m}$ でパラメータ化されたコスト $\transp{\hD}(\mathbf{A}\p-\q)$ を負うものとする。
この緩和の最適値はコストベクトル $\hD$ に直接依存し、次のように書くことができる：
$$\LagrangeMKD(\hD)\eqdef \umin{\p \in \RR^{nm}_+} \transp{\cc}\p-\transp{\hD}(\mathbf{A}\p-\q).$$
まず、この緩和された問題は $\p$ に対する周辺制約を持たないことに注意せよ。この最小化はより多くの $\p$ の解を許すため、$\LagrangeMKD(\hD)$ は $\bar{z}=\MKD_{\C}(\a,\b)$ より小さいと予想される。実際、主問題 \eqref{eq-mk-discr-algo} の任意の最適解を $\p^\star$ と書くと、次が得られる：
$$\umin{\p \in \RR^{nm}_+} \transp{\cc}\p-\transp{\hD}(\mathbf{A}\p-\q) \leq \transp{\cc}\p^\star-\transp{\hD}(\mathbf{A}\p^\star-\q)=\transp{\cc}\p^\star=\bar{z}.$$
したがって上記のアプローチは、任意のコストベクトル $\hD$ に対して元の問題~\eqref{eq-mk-discr-algo}の最適上界を計算するために使用できる問題を定義する。この関数は $\MKD$ の Lagrange 双対関数と呼ばれる。
双対性理論の目標は、\emph{任意の}コストベクトル $\hD$ について $\LagrangeMKD$ を\emph{最大化}することにより、最良の下界 $\underline{z}$ を計算することである。すなわち
$$
	\underline{z}=\umax{\hD} \pa{ \LagrangeMKD(\hD)=\umax{\hD} \transp{\hD}\q + \umin{\p \in \RR^{nm}_+} \transp{(\cc-\transp{A}\hD)}\p }.
$$
$\p$ に関する最小化を含む第2項は、$\transp{\cc}- \transp{A}\hD$ のいずれかの座標が負であれば $-\infty$ となることが容易に示される。
実際、例えばある添字 $i\leq n+m$ に対して $\cc_i-(\transp{A}\hD)_i<0$ であれば、$\p$ として標準基底ベクトル $\e_i$ に任意の大きな正の値を掛けたものを取れば、非有界な値が得られる。
下界 $\LagrangeMKD(\hD)$ を最大化しようとする際、ベクトル $\hD$ を $\transp{A}\hD\leq \cc$ となるものに制限することは理にかなっており、その場合、最良の下界は $$\underline{z}=	\umax{\substack{\hD \in \RR^{n+m}\\\transp{A}\hD \leq \cc}}\transp{\hD}\q$$ となる。したがって、弱双対性の結果、すなわち $\underline{z}\leq \bar{z}$ が証明された。
\end{rem}


\section{$\C$-変換}\label{sec-c-transforms} 本節では、双対 optimal transport 問題~\eqref{eq-mk-discr-dual}の重要な性質を提示する。この性質は、第~\ref{s-c-transform}節の半離散 optimal transport 問題に用いる際により重要な意味を持つ。本節は、行和制約と列和制約に従って双対変数を分割する元の定式化~\eqref{eq-dual}に基づいている：
\eql{\label{eq-dual-discrete-split}\MKD_\C(\a,\b) = \umax{(\fD,\gD) \in \PotentialsD(\C)} \dotp{\fD}{\a} + \dotp{\gD}{\b}.
}
任意の双対実行可能ペア $(\fD,\gD)$ を考える。$\fD$ の値を「凍結」すると、$\gD$ に対する最良のベクトル解は $\fD$ の $\C$-変換ベクトル、$\fD^{\,\C}\in\RR^m$ として定義される以下のものであることに気づく：
\eq{(\fD^{\,\C})_j = \min_{i\in \range{n}} \C_{ij}-\fD_i,}
実際、$(\fD,\fD^{\,\C})\in\PotentialsD(\C)$ であること、および $\fD^{\,\C}$ がこの制約を満たす最大のベクトルであることは容易に証明できる。したがって次が成り立つ：
\eq{\dotp{\fD}{\a} + \dotp{\gD}{\b} \leq \dotp{\fD}{\a} + \dotp{\fD^{\,\C}}{\b}.}
この結果により、まず双対問題を単一の変数 $\fD$ で表された区分的アフィン凹最大化問題として再定式化できる：
\eql{\label{eq-semidual-discret}\MKD_\C(\a,\b) = \umax{\fD\in\RR^{n}} \dotp{\fD}{\a} + \dotp{\fD^{\,\C}}{\b}.}

この結果を脇に置くと、同じ推論は、$\gD$ の値を「凍結」して $\gD$ の $\bar{\C}$-変換、すなわちベクトル $\gD^{\bar{\C}}\in\RR^n$ を考える場合にも当然適用される。これは次のように定義される：
\eq{(\gD^{\bar{\C}})_i = \min_{j\in \range{m}} \C_{ij}-\gD_j,}
目的関数の異なる増加は
\eq{\dotp{\fD}{\a} + \dotp{\gD}{\b} \leq \dotp{\gD^{\bar{\C}}}{\a} + \dotp{\gD}{\b}.}
である。与えられた $\fD$ から出発して、$\fD$ を改善するために $\C$ 変換と $\bar{\C}$ 変換を複数回交互に適用することは魅力的である。実際、次の不等式の列が得られる：
\eq{
	\dotp{\fD}{\a} + \dotp{\fD^{\,\C}}{\b} \leq \dotp{\fD^{\,\C\bar{\C}}}{\a} + \dotp{\fD^{\,\C}}{\b}\leq  \dotp{\fD^{\,\C\bar{\C}}}{\a} + \dotp{\fD^{\,\C\bar{\C}\C}}{\b}\leq  \dots
}
各反復で目的関数が真に増加することを期待するかもしれない。しかし、$\C$ 変換と $\bar{\C}$ 変換の交互適用はすぐにプラトーに達するため、これはうまくいかない。

\begin{prop}\label{prop-ccc-2}以下の等式が成り立つ。ここでベクトル間の不等号は要素ごとに理解される：
	\begin{enumerate}[label=(\roman*)]
	\item $\fD\leq \fD\,' \Rightarrow \fD^{\,\C}\geq \fD\,'^{\,\C}$,
	\item $\fD^{\,\C\bar{\C}} \geq \fD$, $\gD^{\bar{\C}\C} \geq \gD$,
	\item $\fD^{\,\C\bar{\C}\C}=\fD^{\,\C}.$
\end{enumerate}
\end{prop}
\begin{proof} 第一の不等式は $\C$-変換の定義から従う。$\fD^{\,\C\bar{\C}}$ の定義を展開すると
\eq{\left(\fD^{\,\C\bar{\C}}\right)_i= \min_{j\in \range{m}} \C_{ij}-\fD^{\,\C}_j = \min_{j\in \range{m}} \C_{ij}- \min_{i'\in \range{n}} \C_{i'j}-\fD_{i
'}.}
ここで $-\min_{i'\in \range{n}} \C_{i'j}-\fD_{i'} \geq -(\C_{ij}-\fD_i)$ であるから、
\eq{\left(\fD^{\,\C\bar{\C}}\right)_i \geq \min_{j\in \range{m}} \C_{ij}- \C_{ij}+\fD_i = \fD_i.}
が復元される。$\gD^{\bar{\C}\C} \geq \gD$ の関係も同様に得られる。ここで $\gD=\fD^{\,\C}$ とおく。すると $\gD^{\bar{\C}}=\fD^{\,\C\bar{\C}}\geq \fD$ である。したがって、結果 (i) を用いると $\fD^{\,\C\bar{\C}\C}\leq \fD^{\,\C}$ が得られる。結果 (ii) より $\fD^{\,\C\bar{\C}\C}\geq \fD^{\,\C}$ であり、等式が証明される。
\end{proof}


\section{相補スラック性}\label{s-complementary}

主問題~\eqref{eq-mk-discr-primal}と双対問題~\eqref{eq-mk-discr-dual},~\eqref{eq-dual}は、最適な主変数 $\P^{\star}$ と双対変数 $(\fD^{\star},\gD^{\star})$ を得るために独立に解くことができる。次の命題はそれらの関係を特徴づける。

\begin{prop}\label{prop-primal-dual-optimal} $\P^\star$ と $\fD^\star,\gD^\star$ をそれぞれ主問題~\eqref{eq-dual-generic}と双対問題~\eqref{eq-mk-discr}の最適解とする。このとき、任意のペア $(i,j)\in\range{n}\times\range{m}$ に対して $\P^\star_{i,j}(\C_{i,j}-\fD^\star_i+\gD^\star_j)=0$ が成り立つ。言い換えれば、$\P^\star_{i,j}>0$ ならば必然的に $\fD^\star_i+\gD^\star_j=\C_{i,j}$ であり、$\fD^\star_i+\gD^\star_j<\C_{i,j}$ ならば必然的に $\P^\star_{i,j}=0$ である。
\end{prop}
\begin{proof} 強双対性より $\dotp{\P^\star}{\C} = \dotp{\fD^\star}{\a}+\dotp{\gD^\star}{\b}$ が成り立つ。$\P^\star\ones_m = \a$ および $\transp{\P^\star}\ones_n = \b$ を思い出すと、
	$$\begin{aligned}\dotp{\fD^\star}{\a}+\dotp{\gD^\star}{\b} &= \dotp{\fD^\star}{\P^\star\ones_m}+\dotp{\gD^\star}{\transp{\P^\star}\ones_n}\\&=  \dotp{\fD^\star\transp{\ones_m}}{\P^\star}+\dotp{\ones_n\transp{\gD^\star}}{\P^\star},\end{aligned}$$
	となり、$$\dotp{\P^\star}{\C- \fD^\star\oplus \gD^\star}=0$$ が得られる。
$(\fD^\star,\gD^\star)$ は双対制約の多面体~\eqref{eq-feasible-potential}に属するため、行列 $\C- \fD^\star\oplus \gD^\star$ の各成分は必然的に非負である。したがって、$\P$ のすべての成分が非負であることから、上記の内積が $0$ に等しいという制約は、$\P_{i,j}>0$ であるすべての添字のペア $(i,j)$ に対して $\C_{i,j}-(\fD_i+\gD_j)$ がゼロでなければならず、$\C_{i,j}>\fD_i+\gD_j$ であるすべての添字のペア $(i,j)$ に対して $\P_{i,j}=0$ でなければならないことを強制する。
\end{proof}

逆の結果も真である。まず、主問題と双対問題の2つの変数が相補的であるという概念を定義する。

\begin{defn}\label{def-complementary}
行列 $\P\in\RR^{n\times m}$ とベクトルのペア $(\fD,\gD)$ が $\C$ に関して相補的であるとは、$\P_{i,j}>0$ であるすべての添字のペア $(i,j)$ に対して $\C_{i,j}=\fD_i+\gD_j$ も成り立つことをいう。
\end{defn}

実行可能な主変数と双対変数のペアが相補的であれば、それらが最適であると結論できる。

\begin{prop}\label{prop-primal-dual-optimality} $\P$ と $(\fD,\gD)$ がそれぞれ主問題~\eqref{eq-dual-generic}と双対問題~\eqref{eq-mk-discr}の相補的かつ実行可能な解であるとする。このとき、$\P$ と $(\fD,\gD)$ はそれぞれ主最適および双対最適である。
\end{prop}
\begin{proof}
	弱双対性より、
$$\MKD_{\C}(\a,\b)\leq \dotp{\P}{\C}=\dotp{\P}{\fD\oplus\gD}=\dotp{\a}{\fD}+\dotp{\b}{\gD}\leq \MKD_{\C}(\a,\b)$$
が成り立ち、したがって $\P$ と $(\fD,\gD)$ はそれぞれ主最適および双対最適である。
\end{proof}


\section{輸送多面体の頂点}\label{s-extremal}

凸集合の頂点または端点とは、その集合の点 $\mathbf{x}$ であって、$\mathbf{x}=(\mathbf{y}+\mathbf{z})/2$ となるその集合の $\mathbf{y}$ と $\mathbf{z}$ が存在するならば必然的に $\mathbf{x}=\mathbf{y}=\mathbf{z}$ となるものであることを思い出そう。空でない有界な実行可能集合を持つ線形計画は、実行可能集合の頂点（または端点）で最小値を達成する~\citep[p.~65, Theo.~2.7]{bertsimas1997introduction}。主 optimal transport 問題~\eqref{eq-mk-discr-primal}の実行可能集合 $\CouplingsD(\a,\b)$ は有界であるから、最適な $\P$ の探索を多面体 $\CouplingsD(\a,\b)$ の端点の集合に限定できる。$\CouplingsD(\a,\b)$ において端点的な行列 $\P$ は、広範な研究の対象となってきた興味深い構造を持つ~\citep[\S8]{brualdi2006combinatorial}。その構造を記述するには、二部グラフの形式で輸送問題を表す必要がある。

\subsection{$\CouplingsD(\a,\b)$ のすべての頂点の台の木構造}
$V=(1,2,\dots,n)$ と $V'=(1',2',\dots,m')$ を2つのノード集合とする。集合 $V'$ のラベルには $V$ のものと区別するためにプライムを付ける。それらの和集合 $V\cup V'$（$n+m$ 個のノードを持つ）と、それらの間のすべての $nm$ 本の有向辺の集合 $\alledges = \{ (i,j'), i \in \range{n}, j\in \range{m}\}$ を考える（ここでは $V'$ の $j'$ を形成するために整数 $j\leq m$ にプライムを付けるだけである）。各辺 $(i,j')$ には対応するコスト値 $\C_{ij}$ を関連付ける。$V$ と $V'$ の間の完全二部グラフ $\mathcal{G}$ は $(V\cup V',E)$ である。輸送計画はこのグラフ上のフローであり、図~\ref{fig-simplex}に非形式的に示されるように、ソース制約（各ノード $i$ から $\a_i$ が流出）とシンク制約（各ノード $j'$ に $\b_j$ が流入）を満たす。$\CouplingsD(\a,\b)$ の端点は以下の性質を持つ~\citep[p.~338,~Theo.~8.1.2]{brualdi2006combinatorial}。

\begin{figure}[h!]
\centering
\includegraphics[width=\linewidth]{algo-basics-tree-solution/networksol}
\caption{\label{fig-simplex}
二部ネットワークフロー問題としての optimal transport 問題。ここで $n=3,m=4$ である。ソースヒストグラム $\a$ のすべての座標は、左側のソースノードとして $1,2,3$ とラベル付けされ、ターゲットヒストグラム $\b$ のすべての座標はノード $1',2',3',4'$ としてラベル付けされる。グラフは、すべてのソースノードがすべてのターゲットノードに接続され、追加の辺がないという意味で二部的である。各辺 $(i,j')$ にはコスト $\C_{ij}$ が関連付けられる。右側に実行可能なフローが示されている。命題~\ref{prop-extremal}は、このフローが少なくとも1つの巡回 $((1,1'),(2,1'),(2,4'),(1,4'))$ を持つため端点的でないことを示す。
}
\end{figure}


\begin{figure}[h!]
\centering
\includegraphics[width=\linewidth]{algo-basics-tree-solution/tree-sol2}
\caption{\label{fig-perturb}
台のグラフに巡回を持つ解 $\P$ は、$\P$ がそれらの平均となるような2つの実行可能な解 $\Q$ と $\RRd$ を得るために摂動でき、したがって $\P$ が端点的であることが否定される。
}
\end{figure}

\begin{prop}[端点的解]\label{prop-extremal}
$\P$ を多面体 $\CouplingsD(\a,\b)$ の端点とする。$S(\P)\subset \alledges$ を辺の部分集合 $\{(i,j'), i\in\range{n}, j\in\range{m} \text{ such that } \P_{ij}>0\}$ とする。このとき、グラフ $G(\P)\eqdef (V\cup V',S(\P))$ は巡回を持たない。特に、$\P$ は $n+m-1$ 個を超える非零成分を持つことはできない。
 \end{prop}
\begin{proof}
背理法により証明する。$\P$ が多面体 $\CouplingsD(\a,\b)$ の端点であり、その対応する辺の集合 $S(\P)$（簡潔に $F$ と記す）が、グラフ $G=(V\cup V',F)$ が巡回を含むようなものであると仮定する。すなわち、$k>1$ と相異なる添字の列 $i_1,\dots,i_{k-1}\in\range{n}$ および $j_1,\dots,j_{k-1}\in\range{m}$ が存在して、以下の辺の集合 $H$ が $F$ の部分集合をなす。$$H=\left\{(i_1,j_1'), (i_2,j_1'), (i_2,j_2'),\dots,(i_k,j_k'),(i_1,j_k')\right\}.$$
	ここで、$\P=(\Q+\RRd)/2$ となる2つの実行可能な行列 $\Q$ と $\RRd$ を構成する。そのために、$H$ に対応する\emph{有向}巡回 $\bar{H}$、すなわち対の列 $i_1\rightarrow j_1', j_1' \rightarrow i_2, i_2 \rightarrow j_2',\dots, i_k \rightarrow j_k' , j_k' \rightarrow i_1$、および基本的なフロー量 $\varepsilon < \min_{(i,j')\in F}\P_{ij}$ を考える。
	摂動行列 $\E$ を、$(i,j)$ 成分が $i\rightarrow j' \in \bar{H}$ ならば $\varepsilon$、$j\rightarrow i'\in \bar{H}$ ならば $-\varepsilon$、それ以外は零であるものとして定義する。図~\ref{fig-perturb}に示されるように、行列 $\Q=\P+\E$ と $\RRd=\P-\E$ を定義する。$\varepsilon$ が十分に小さいため、$\Q$ と $\RRd$ のすべての要素は非負である。構成により、$\E$ は、すべての成分が $0$ の行（もしくは列）を持つか、$i_1,\dots,i_k$（もしくは $j_1,\dots,j_k$）で添字付けられたものについて正確に1つの成分が $\varepsilon$ でもう1つが $-\varepsilon$ であるかのいずれかである。したがって、$\E$ は $\E\ones_{m}=\zeros_{n}$ かつ $\transp{\E}\ones_{n}=\zeros_{m}$ を満たし、$\Q$ と $\RRd$ は $\P$ と同じ周辺を持ち、したがって実行可能である。最後に $\P=(\Q+\RRd)/2$ であり、$\Q,\RRd\ne \P$ であるから、$\P$ が端点であるという事実に矛盾する。$k$ 個のノードを持つ巡回のないグラフは $k-1$ 本を超える辺を持てないから、$S(\P)$ は $n+m-1$ 本を超える辺を持てず、したがって $\P$ は $n+m-1$ 個を超える非零成分を持つことはできない。
\end{proof}

\subsection{北西隅規則}\label{subsec-northwest}
北西（NW）隅規則は、最大 $n+m$ 回の操作で多面体 $\CouplingsD(\a,\b)$ の頂点を生成するヒューリスティックである。このヒューリスティックは、次節で概説するネットワークシンプレックスのような主変数に作用する任意のアルゴリズムの初期化に役立ちうる。

この規則は、$\P_{1,1}$ に $\min(\a_1,\b_1)$ を設定して可能な限り大きな値を与えることから始まる。各ステップにおいて、成分 $\P_{i,j}$ は行制約 $i$、列制約 $j$、または可能であればその両方を飽和するように選ばれる。その後、添字 $i,j$ は次のように更新される：最初の場合は $i$ が増加し、2番目の場合は $j$ が増加し、3番目の場合は $i$ と $j$ の両方が増加する。この規則は $\P_{n,m}$ が値を受け取るまで進む。

形式的には、アルゴリズムは次のように動作する：$i$ と $j$ を $1$ に、$r\leftarrow\a_1,c\leftarrow \b_1$ に初期化する。$i\leq n$ かつ $j\leq m$ の間、$t\leftarrow \min(r,c)$, $\P_{i,j}\leftarrow t$, $r\leftarrow r-t$, $c\leftarrow s-t$ とおく。$r=0$ ならば $i$ を増加させ、$i\leq n$ ならば $r\leftarrow\a_i$ と更新する。$c=0$ ならば $j$ を増加させ、$j\leq n$ ならば $c\leftarrow \b_j$ と更新する。これを繰り返す。$\a=[0.2,0.5,0.3]$ かつ $\b=[0.5,0.1,0.4]$ を仮定した場合のこの列の例を示す：
$$\begin{aligned}\begin{bmatrix} \bullet & 0 & 0 \\ 0 & 0 & 0 \\ 0& 0 & 0\end{bmatrix} &\rightarrow \begin{bmatrix} 0.2 & 0 & 0 \\ \bullet & 0 & 0 \\ 0& 0 & 0\end{bmatrix} &\rightarrow \begin{bmatrix} 0.2 & 0 & 0 \\ 0.3 & \bullet & 0 \\ 0& 0 & 0\end{bmatrix}\\ &\rightarrow \begin{bmatrix} 0.2 & 0 & 0 \\ 0.3 &0.1 &\bullet \\ 0& 0 & 0\end{bmatrix} &\rightarrow \begin{bmatrix} 0.2 & 0 & 0 \\ 0.3 &0.1 &0.1 \\ 0& 0 & \bullet\end{bmatrix} & \rightarrow \begin{bmatrix} 0.2 & 0 & 0 \\ 0.3 &0.1 &0.1 \\ 0& 0 & 0.3\end{bmatrix}\end{aligned}$$
このヒューリスティックを通じて得られる一意の計画を $\NW(\a,\b)$ と書く。

ただし、$\a$ と $\b$ の順序を任意に置換し、対応する NW 隅テーブルを計算し、列と行の順序を再び反転して $\CouplingsD(\a,\b)$ のテーブルを復元することで得られる、はるかに多数の NW 隅解が存在することに注意せよ。$\sigma=(3,1,2),\sigma'=(3,2,1)$ とおくと $\a_\sigma=[0.3,0.2,0.5], \b_{\sigma'}=[0.4,0.1,0.5]$ となり、$\sigma^{-1}=(2,3,1),\sigma'=(3,2,1)$ である。次が成り立つ：
\begin{gather*}
\NW(\a_\sigma,\b_{\sigma'}) = \begin{bmatrix} 0.3 & 0 & 0 \\ 0.1 & 0.1 & 0 \\ 0& 0 & 0.5\end{bmatrix} \in \CouplingsD(\a_\sigma,\b_{\sigma'}),\\
\NW_{\sigma^{-1}\sigma'^{-1}}(\a_\sigma,\b_{\sigma'})= \begin{bmatrix} 0 & 0.1 & 0.1 \\ 0.5 & 0 & 0 \\ 0& 0 & 0.3\end{bmatrix}\in \CouplingsD(\a,\b).
\end{gather*}

$\mathcal{N}(\a,\b)$ をこの方法で生成できるすべての NW 隅解の集合とする：
$$\mathcal{N}(\a,\b)\eqdef\{ \NW_{\sigma^{-1}\sigma'^{-1}}(r_\sigma,c_{\sigma'}), \sigma,\sigma'\in S_d\}.$$
すべての NW 隅解は構成上、最大 $n+m-1$ 個の非零要素を持つ。NW 隅規則は $\a_\sigma$ と $\b_\sigma'$ に対して構成上一意のテーブルを生成するが、同じテーブルを生成しうる行/列の置換のペア $(\sigma,\sigma')$ は指数的に多く存在する~\citep[p. 2]{stougie2002polynomial}。$\mathcal{N}(\a,\b)$ は $\CouplingsD(\a,\b)$ の端点の集合の（通常は真の）部分集合をなす~\citep[Cor. 8.1.4]{brualdi2006combinatorial}。



\section{ネットワークシンプレックスのヒューリスティックな説明}\label{s-networksimplex}

グラフ $G(\P)=(V\cup V',S(\P))$ が巡回を持たない実行可能な行列 $\P$ を考える。$\P$ は $n+m-1$ 個以下の非零成分を持ち、命題~\ref{prop-extremal}により $\CouplingsD(\a,\b)$ の頂点である。
命題~\ref{prop-primal-dual-optimality}に従い、$\P$ に対して実行可能（\ie $\C-\fD\oplus\gD$ が非負の成分を持つ）かつ相補的（$S(\P)$ 内の添字のペア $(i,j')$ が $\C_{i,j}=\fD_i+\gD_j$ を満たす）な双対解 $(\fD,\gD)$ を得れば、$\P$ の最適性を証明するのに十分である。
ネットワークシンプレックスは2つの単純な原理に依存する：各実行可能な主解 $\P$ に相補的なペア $(\fD,\gD)$ を関連付けることができる。
そのペアが実行可能であれば、最適性に到達したことになる。そうでない場合、すなわち $\fD_i+\gD_j> \C_{i,j}$ となる $i,j$ が存在する場合、$\P$ の修正を考えることができる。その修正は実行可能性を維持しつつ、相補的なペア $(\fD,\gD)$ が実行可能性に近づくように変更される。

\subsection{$\P$ に相補的な双対ペアの取得}\label{subsec-obtaining}
シンプレックスは、まず任意の端点的解 $\P$ に相補的な双対変数のペア $(\fD,\gD)$ を関連付けることから進む。これは、$S(\P)$ 内の任意の $(i,j')$ に対して $\fD_i+\gD_j$ が $\C_{i,j}$ に等しくなるような2つのベクトル $\fD$ と $\gD$ を見つけることによって単純に実行される。ただし、これ自体は $(\fD,\gD)$ が実行可能であることを保証しないことに注意せよ。

\begin{figure}[h!]
	\centering
	\includegraphics[width=\linewidth]{algo-basics-tree-solution/f_of_P}
	\caption{\label{fig-f_of_P}
	実行可能な輸送 $\P$ とそれに対応する辺の集合 $S(\P)$ およびグラフ $G(\P)$。見られるように、グラフ $G(\P)=(\{1,\dots,5,1',\dots,6'\},S(\P))$ は森、すなわち木グラフの和集合として表現でき、この場合は3つの木からなる。}
\end{figure}


$s$ を $S(\P)$ の濃度とする。$\P$ は端点的であるから $s\leq n+m-1$ である。$G(\P)$ は巡回を持たないから、図~\ref{fig-f_of_P}に示されるように、$G(\P)$ は木または森（木の和集合）のいずれかである。$\P$ に相補的なペア $(\fD,\gD)$ を目指して、$n+m$ 個の変数に対する以下の $s$ 個の線形等式制約を考える：
\begin{equation}\label{eq-dual-variables}
	\begin{array}{ccc}
\fD_{i_1}+\gD_{j_1}&=&\C_{i_1,j_1}\\
\fD_{i_2}+\gD_{j_1}&=&\C_{i_2,j_1}\\
\vdots & = & \vdots\\
\fD_{i_s}+\gD_{j_s}&=&\C_{i_s,j_s},\\
\end{array}
\end{equation}
ここで $S(\P)$ の要素は $(i_1,j_1'),\dots,(i_s,j_s')$ と列挙される。

$s\leq n+m-1 < n+m$ であるから、上記の線形系~\eqref{eq-dual-variables}は常に不定である。この退化は、$n+m$ 個の制約による $\CouplingsD(\a,\b)$ のパラメータ化が $n+m$ 個の双対変数をもたらすことに部分的に起因すると解釈できる。注意~\ref{rem-transportation-polytope}で概説されたより慎重な定式化では、$n+m-1$ 個の制約のみを持つ等価な定式化、したがって $n+m-1$ 個の双対変数が得られたであろう。しかし、$s$ は $n+m-1$ より真に小さくなることもありうる。これは $G(\P)$ が2つ以上の木の非連結和である場合に起こる。例えば、図~\ref{fig-f_of_P}では $5+6=11$ 個の双対変数（各ノードに1つ）があるが、これら $11$ 個のノードの間には $8$ 本の辺しかなく、すなわち $(\fD,\gD)$ を定義するための $8$ 個の線形方程式しかない。したがって、この設定下で不定な双対変数の数は $G(\P)$ の連結成分の数と同じになる。

$G(\P)$ に列挙された木の中の1つを考えよう。その木がソースノードの中に $k$ 個のノード $i_1,\dots,i_k$、ターゲットノードの中に $l$ 個のノード $j_1',\dots,j_{l}'$ を持ち、$r\eqdef k+l$ 個のノードと $r-1$ 本の辺を持つとする。これは $\fD$ の $k$ 個の変数と $\gD$ の $l$ 個の変数に対応し、$r-1$ 個の線形方程式で結ばれている。不定性を解消するために、その木の中の根ノードを任意に選び、対応する双対変数に値 $0$ を割り当てることができる。そこから、幅優先探索または深さ優先探索を用いて木を走査し、図~\ref{fig-tree-linear-solve}に示されるように、その木の他のすべての双対変数の値を決定する単純な変数代入の列を得ることができる。この手続きを $\P$ のグラフのすべての木に対して繰り返すことで、$\P$ に相補的な双対変数のペア $(\fD,\gD)$ が得られる。

\begin{figure}[h!]
	\centering
	\includegraphics[width=\linewidth]{algo-basics-tree-solution/tree-linear-solve}
	\caption{\label{fig-tree-linear-solve}
 図~\ref{fig-f_of_P}に示されたグラフ $G(\P)$ の最初の木に現れる5つのノードに対応する5つの双対変数 $\fD_1,\fD_2,\gD_1,\gD_2,\gD_3$ は、コスト行列 $\C$ の対応する成分を含む4つの線形方程式で結ばれている。この系は退化しているため、その木の根（この例ではノード 1）を選び、対応する変数を $0$ に設定し、根から木を走査（幅優先探索または深さ優先探索のいずれか）して、残りの4つの双対変数の値を反復的に得る。}
\end{figure}

\subsection{ネットワークシンプレックスの更新}
前節で得られた双対ペア $(\fD,\gD)$ は、すべての $i,j$ に対して $\fD_i+\gD_j\leq \C_{i,j}$ が成り立つという意味で実行可能かもしれず、その場合は命題~\ref{prop-primal-dual-optimality}により最適解に到達したことになる。そうでない場合、すなわち $\fD_i+\gD_j> \C_{i,j}$ となる $i,j$ が存在する場合、ネットワークシンプレックスアルゴリズムが始動する。まず、グラフ $G$ を実行可能な解 $\P$ に対応するグラフ $G(\P)$ に等しく初期化し、違反辺 $(i,j')$ を $G$ に追加する。すると2つの場合が生じうる：
\begin{enumerate}[label={(\alph*)}]
	\item $G$ が（依然として）森である場合。これは $(i,j')$ が2つの既存の部分木を結合する場合に起こりうる。第~\ref{subsec-obtaining}節で概説したアプローチを $G$ に対して用いることで、新しい相補的双対ベクトル $(\fD,\gD)$ を回復できる。この追加は $n+m$ 個の双対変数の不定性を単に除去するだけであり、主変数 $\P$ の変更は生じないことに注意せよ。この更新は通常、退化的と呼ばれる。$(i,j')$ がグラフ $G$ に入ったにもかかわらず $\P_{i,j}$ は $0$ のままだからである。ただし、$G(\P)$ は $G$ に含まれる。

	\item $G$ が巡回を持つ場合。この場合、$G$ が依然として森であることを保証するために $G$ から辺を除去する必要があるが、同時に $\P$ を修正して $\P$ が実行可能であり $G(\P)$ が $G$ に含まれ続けるようにしなければならない。これらの操作はすべて、$\P_{i,j}$ の値を増加させ、検出された巡回に現れる $\P$ の他の成分を、命題~\ref{prop-extremal}の証明で用いたものと非常に類似した方法で修正することにより実行できる。より正確には、その巡回を $(i_1,j_1'), (j_1',i_2), (i_2,j_2'),\dots, (i_l,j_l'),(j_l',i_{l+1})$ と書く。ここで、経路が $i$ で始まり $i$ で終わる巡回であることを保証するために $i_1=i_{l+1}=i$ とし、巡回が追加された辺 $\{i,j\}$ から始まることを強調するために $j_1=j$ とする。すべての「正の」辺 $(i_k,j_k')$（$k\leq l$）のフローを増加させ、「負の」辺 $(j_k',i_{k+1})$（$k\leq l$）のフローを減少させて、更新された主解 $\tilde{\P}$ を得る。これは以下の成分を除いて $\P$ に等しい：
$$\forall k\leq l,\quad \tilde{\P}_{i_k,j_k}:= \P_{i_k,j_k}+\theta;\quad \tilde{\P}_{i_{k+1},j_k}:= \P_{i_{k+1},j_k}-\theta.$$
ここで $\theta$ はその巡回を用いた添字 $i,j$ での最大可能な増加量である。$\theta$ の値は、巡回によって負の影響を受ける最小のフロー、すなわち $\min_{k} \P_{i_{k+1},j_k}$ によって制御される。この更新は図~\ref{fig-simplex-update}に示されている。$k^\star$ をこの最小を達成する添字とする。その後、$(i_{k^\star+1},j_{k^\star})$ を $G$ から除去して更新を閉じ、第~\ref{subsec-obtaining}節で概説したアプローチを用いて新しい双対変数 $(\fD,\gD)$ を計算する。
\end{enumerate}

\subsection{主解の改善} これは必ずしも我々の当初の動機ではなかったが、上記の操作は $\P$ のコストを改善することしかできないことを示すことができる。追加された辺が巡回を作成しなかった場合、すなわち上記の場合 {(a)} では、主解は変化しない。巡回が作成された場合、すなわち場合 {(b)} では、$\P$ は $\tilde{\P}$ に更新され、次の等式が成り立つ：
\eq{
\dotp{\tilde{\P}}{\C}-\dotp{\P}{\C}= \theta \left(\;\sum_{k=1}^{l} \C_{i_k,j_k} - \sum_{k=1}^{l}\C_{i_{k+1},j_k}\right).
}
ここで、前の反復の最後に計算された双対ベクトル $(\fD,\gD)$ を用いる。それらは、$G$ に元からあった\emph{すべての}辺に対して $\f_{i_k}+\g_{i_k}=\C_{i_k,j_k}$ かつ $\f_{i_{k+1}}+\g_{i_k}=\C_{i_{k+1},j_k}$ を満たし、次の等式が得られる：
\eq{
\begin{aligned}\sum_{k=1}^{l} \C_{i_k,j_k} - \sum_{k=1}^{l}\C_{i_{k+1},j_k} &= \C_{i,j} + \sum_{k=2}^{l} \fD_{i_k}+\gD_{j_k} - \sum_{k=1}^{l} \fD_{i_{k+1}}+\gD_{j_k}\\
	&= \C_{i,j} - (\fD_{i} +\gD_{j}).
\end{aligned}
}
この項は定義により負である。$i,j$ は $C_{i,j} < \fD_{i} -\gD_{j}$ であるために選ばれたからである。したがって、$\theta>0$ ならば \eq{\dotp{\tilde{\P}}{\C}=\dotp{\P}{\C}+\theta\left(\C_{i,j}-(\fD_{i} -\fD_{g})\right)<\dotp{\P}{\C}.}
が成り立つ。$\theta=0$ の場合（$G$ と $G(\P)$ が異なる場合に起こりうる）、グラフ $G$ は単に変更されるが、$\P$ は変更されない。


\begin{figure}[h!]
	\centering
	\includegraphics[width=\linewidth]{algo-basics-tree-solution/simplex-update}
	\caption{\label{fig-simplex-update}
グラフ $G(\P)$ に辺 $\{i,j\}$ を追加すると、(a) この追加後もグラフが森のままである場合、第~\ref{subsec-obtaining}節で概説したアプローチに従って $\fD,\gD$ を再計算できる。(b.1) その辺の追加が巡回を作成する場合、そこから有向経路を定義できる。(b.2) その経路を用いて $\P_{i,j}$ の値を増加させ、フロー実行可能性制約を維持するために巡回に沿ってその変化を伝播させることができる。巡回によって負の影響を受ける辺のうちの1つのフローが $0$ に減少するまでこれを行う。これにより巡回が除去され、$\P$ が更新される。}
\end{figure}

ネットワークシンプレックスアルゴリズムは次のように要約できる：第~\ref{subsec-northwest}節で扱った NW 隅規則などにより与えられる端点的解 $\P$ でアルゴリズムを初期化する。グラフ $G$ を $G(\P)$ で初期化する。第~\ref{subsec-obtaining}節で記述した $G$ の木構造を用いた線形系の解法により、$\P$ に相補的な双対変数のペア $(\fD,\gD)$ を計算する。{(i)} 制約 $\C- \fD\oplus \gD\geq 0$ に対する違反ペアの添字を探す。なければ $\P$ は最適であり停止する。違反ペア $(i,j')$ が存在すれば、{(ii)} 辺 $(i,j')$ を $G$ に追加する。$G$ が依然として巡回を持たなければ、$(\fD,\gD)$ をそれに応じて更新する。巡回があれば、$(i,j')$ が正とラベル付けされるようにそれを有向化し、最小のフロー値を持つ負の辺をその巡回から除去し、図~\ref{fig-simplex-update}に示されるように $\P,G$ を更新し、それに応じて相補的ペア $\fD,\gD$ を構築する。{(i)} に戻る。上記の操作の一部はグラフ操作（巡回検出、木の走査）を必要とし、これらはこの文脈において（\cite[\S5]{bertsekas1998network}）に記述されるように効率的に実装できる。

\citet{Orlin1997}はネットワークシンプレックスの多項式時間計算量を初めて証明した。\citet{Tarjan1997}はその直後に、ピボット辺の選択を支援するより効率的なデータ構造に依存する改善された上界 $O\left(\,(n+m)nm \log(n+m)\log\left((n+m)\|\C\|_\infty\right)\,\right)$ を提供した。



\section{双対上昇法}\label{s-dual-ascent}

双対上昇法はネットワークシンプレックスより数十年前に遡り、\citet{borchardt1865investigando}の研究、およびその後の K\"onig と Egerv\'ary の研究にまで遡ることができる。これは\citet{Kuhn1955}によって回顧されている。ハンガリアン法はこの族で最もよく知られたアルゴリズムであり、$\a$ と $\b$ が等しくともに一様、すなわち $\a=\b=\ones_n/n$ である特殊な場合にのみ機能する。以下では、より一般的な双対上昇法の族の簡潔な説明を与える。これには最大フロー問題の知識が必要である（\cite[\S7.5]{bertsimas1997introduction}）。上記で主問題において提示されたネットワークシンプレックスとは対照的に、双対上昇法は各反復で双対実行可能な解を維持し、$\fD$ と $\gD$ にスパースなベクトルを加えることでその目的関数を漸進的に改善する。我々の説明は主に（\cite[\S7.7]{bertsimas1997introduction}）に由来し、以下の定義から始まる。

\begin{defn} $S\subset\range{n},S'\subset\range{m}'\eqdef \{1',\dots,m'\}$ に対して、$\ones_S$ は $S$ に列挙された添字以外がすべて零で添字の位置が $1$ である $\RR^n$ のベクトルを表し、同様に $\ones_{S'}$ は $S'$ の添字を持つ $\RR^m$ のベクトルを表す。
\end{defn}

以下では、$(\fD,\gD)$ を $\PotentialsD(\C)$ の実行可能な双対ペアとする。これは単に、すべてのペア $(i,j')\in\range{n}\times\range{m}'$ に対して $\fD_i+\gD_j\leq\C_{ij}$ が成り立つことを意味することを思い出そう。$(i,j')$ が $\fD_i+\gD_j=\C_{ij}$ のとき\emph{均衡}ペア（または辺）であると言い、$\fD_i+\gD_j<\C_{ij}$ のとき\emph{非活性}であると言う。この規約のもと、実行可能な双対ペア $(\fD,\gD)$ が集合 $S$ と $S'$ で添字付けられたスパースベクトルを用いてどのように摂動されても実行可能性を保つことができるかを記述する単純な結果から始める。

\begin{prop}\label{prop-feasibility-dual-pair}
$(\tilde{\fD},\tilde{\gD})\eqdef(\fD,\gD)+\varepsilon(\ones_S,-\ones_{S'})$ は、すべての $i\in S$ に対して $(i,j')$ が均衡であることが $j'\in S'$ を含意する場合、十分小さな $\varepsilon>0$ に対して双対実行可能である。
\end{prop}
\begin{proof}
任意の $i\in S$ に対して、$(i,j')$ が非活性であるような、すなわち $\fD_i+\gD_j<\C_{ij}$ であるようなすべての $j'\in\range{m}'$ の集合 $\mathcal{I}_{i}$ を考える。$\varepsilon_i \eqdef \min_{j\in I_{i}} \C_{i,j}-\fD_i-\gD_j$ を、$j'\in \mathcal{I}_i$ に対応する制約に違反することなく $\fD_i$ を増加させることができる最小のマージンと定義する。実際、$\varepsilon\leq \varepsilon_i$ ならば任意の $j'\in \mathcal{I}_i$ に対して $\tilde{\fD}_i+\tilde{\gD}_j<\C_{i,j}$ が成り立つ。次に $i$ に関連する均衡辺の集合 $\mathcal{B}_i$ を考える。$\mathcal{B}_i=\range{m}'\setminus \mathcal{I}_i$ であることに注意せよ。上記の仮定は $j'\in\mathcal{B}_i\Rightarrow j'\in S'$ である。したがって、$j'\in\mathcal{B}_i$ に対して $\tilde{\fD}_i+\tilde{\gD}_j=\fD_i+\gD_j=\C_{i,j}$ が成り立つ。結果として、不等式 $\tilde{\fD}_i+\tilde{\gD}_j\leq \C_{i,j}$ は任意の $j\in \range{m}'$ に対して保証される。許容される最小の増加量、すなわち $\min_{i\in S} \varepsilon_i$ より小さい増加量 $\varepsilon$ を選ぶと、$(\tilde{\fD},\tilde{\gD})$ が双対実行可能であることが復元される。
\end{proof}

第~\ref{subsec-obtaining}節で提示されたネットワークシンプレックスの反復の主な動機は、実行可能な主解 $\P$ から出発して、相補的な実行可能な双対ペア $(\fD,\gD)$ を得ることであった。その目標に到達するため、$\P$ はその相補的双対ペアが双対実行可能性に到達するように漸進的に修正される。対称的なアプローチ、すなわち実行可能な双対変数から出発して実行可能な主変数 $\P$ を得ることが、双対上昇法の動機である。以下の命題は、$(\fD,\gD)$ に対して実行可能性を維持する上昇方向の存在を（構成的に）保証するという意味で、双対上昇法の主要なエンジンである。その方向は、ネットワークシンプレックスと同様に、候補となる主解 $\P$ を設計し、その非実行可能性が $(\fD,\gD)$ の更新を導くことによって構築される。

\begin{prop}\label{prop-dual-ascent-dir} $(\fD,\gD)$ が問題~\eqref{eq-dual-discrete-split}に対して最適であるか、あるいは $S\subset\range{n},S'\subset\range{m}'$ が存在して $(\tilde{\fD},\tilde{\gD})\eqdef(\fD,\gD)+\varepsilon(\ones_S,-\ones_{S'})$ が十分小さな $\varepsilon>0$ に対して実行可能であり、かつ真に良い目的関数値を持つかのいずれかである。
\end{prop}
\begin{proof}
まず $(\fD,\gD)$ に相補的な主変数 $\P$ を考える。そのために、均衡辺の集合 $\mathcal{B}$、すなわち $\fD_i+\gD_j=\C_{i,j}$ であるようなすべてのペア $(i,j')\in\range{n}\times\range{m}'$ を取り、頂点 $\{1,\dots,n,1',\dots,m'\}$ が $\mathcal{B}$ の辺のみで結ばれた二部グラフを形成する。これにソースノード $s$ を加え、\emph{容量付き}辺ですべてのノード $i\in\range{n}$ にそれぞれ容量 $\a_i$ で接続し、同様に終端ノード $t$ もすべてのノード $j'\in\range{m}'$ にそれぞれ容量 $\b_j$ の辺で接続する。これは図~\ref{fig-dualascent}に示されている。Ford--Fulkerson アルゴリズム（\cite[p. 305]{bertsimas1997introduction}）を用いてそのネットワーク上の最大フロー $\mathbf{F}$ を計算できる。すなわち、$(i,j')\in \mathcal{B}$ で添字付けられた $n+m+|\mathcal{B}|$ 個の非負の値 $f_{si}\leq \a_i, f_{ij'}, f_{j't}\leq\b_j$ の族であり、フロー保存制約を満たし $\sum_i f_{si}$ が最大となるものである。そのフロー $\mathbf{F}$ のスループットが $1$ に等しければ、構成上 $\fD,\gD$ に相補的な実行可能な主解 $\P$ を $\mathbf{F}$ から $\P_{i,j}=f_{ij'}$（$(i,j')\in \mathcal{B}$）およびそれ以外を零として定義することで抽出でき、命題~\ref{prop-primal-dual-optimality}により $(\fD,\gD)$ と $\P$ の最適性が得られる。$\mathbf{F}$ のスループットが $1$ より真に小さい場合、ラベリングアルゴリズムは、$\mathbf{F}$ が容量制約を飽和させない $s$ から反復的に到達可能なノードをラベリング（特定）し、さらにラベル付けされたノードにフローを寄与するノードもラベル付けすることにより進む。ラベル付けされたノードは空でない集合 $Q$ に格納され、$\mathbf{F}$ の最適性により終端ノード $t$ は含まない（厳密な提示は\citealt[p. 308]{bertsimas1997introduction}を参照）。$Q$ は2つの集合 $S=Q\cap\range{n}$ と $S'=Q\cap\range{m}'$ に分割できる。合計スループットが $1$ より真に小さいと仮定したため、$S\ne\emptyset$ である。まず、$i\in S$ かつ $(i,j)$ が均衡であれば、$j'$ は必然的に $S'$ に属することに注意せよ。実際、すべての辺 $(i,j')$ は構成上無限容量を持つため、ラベリングアルゴリズムは $i$ を $S$ に含めれば必然的に $j'$ に到達する。命題~\ref{prop-feasibility-dual-pair}により、$\tilde{\fD},\tilde{\gD}$ の実行可能性を保証するのに十分小さな $\varepsilon$ が存在する。$(\tilde{\fD},\tilde{\gD})$ が $(\fD,\gD)$ より良い目的関数値を持つことを保証するために $\transp{\ones_S}{\a}-\transp{\ones_{S'}}{\b}>0$ を証明する必要がある。$\bar{S}=\range{n}\setminus S$、$\bar{S'}=\range{m}'\setminus S'$ とおき、
$$A = \sum_{i\in S} f_{si}, \quad B = \sum_{i\in \bar{S}} f_{si}, \quad C = \sum_{j'\in S'} f_{j't}, \quad D = \sum_{j'\in \bar{S'}} f_{j't}$$
と定義する。合計最大フローは $s$ から始まるため $A+B$ に等しいが、$t$ に到着するため $C+D$ にも等しい。フロー保存制約により同じフローは $B+C$ にも等しく、したがって $A=C$ である。一方、ラベリングアルゴリズムの定義により $S$ 内のすべての $i$ に対して $f_{si}<\a_i$ が成り立ち、$\bar{S}'$ 内の $j'$ に対しては、考えているフローの最適性により $t$ が $S'$ に入れないため $f_{j't}=\b_j$ である。したがって $A<\transp{\ones_S}\a$ かつ $C=\transp{\ones_S'}\b$ である。ゆえに $\transp{\ones_S}\a-\transp{\ones_S'}\b>A-C=0$ である。
\end{proof}

\begin{figure}[h!]
\centering
\includegraphics[width=\linewidth]{algo-basics-tree-solution/dual_ascent}
\caption{\label{fig-dualascent}
図~\ref{fig-f_of_P}で最初に導入された周辺を含む輸送問題を考える（$n=5, m=6$）。2つの実行可能な双対ベクトル $\fD,\gD$ が与えられたとき、$(\fD,\gD)$ に相補的な「最良の」フロー行列 $P$ を得ようとする。これは、$\fD_i+\gD_j=\C_{i,j}$ となる添字に対応する辺 $(i,j')$ 上でのみ $\P$ が正の値を取りうることを意味する。プロット (a) では点線で表されている。そのグラフ構造で達成可能な最良のフローは、抽象的なソースノード $s$ から $i\in\range{n}$ とラベル付けされたすべてのノードに接続し、抽象的な終端ノード $t$ から $j'\in\range{m'}$ とラベル付けされたすべてのノードに接続する容量付きネットワークにおける最大フロー問題として定式化でき、辺 $(s,i),(j',t)$（$i\in \range{n},j\in\range{m}$）の容量はそれぞれ $\a_i,\b_j$ であり、他はすべて無限容量である。Ford--Fulkerson アルゴリズム（\cite[p. 305]{bertsimas1997introduction}）を適用してそのような最大フローを計算でき、プロット (b) に示されるように、問題を解くために必要な $1$ 単位の質量のうち $0.74$ 単位しか達成できない。最大フローアルゴリズムが使用するサブルーチンの1つであるラベリングアルゴリズム（\cite[p. 308]{bertsimas1997introduction}）を用いて、$s$ から飽和していないフローを受け取るノード（および再帰的にそのすべての後続ノード）を特定でき、プロット (c) ではオレンジの線で示されている。ラベリングアルゴリズムはデフォルトで、ラベル付けされたノードに正のフローを送るノードも追加する。これはノード $3$ を選択するために使用された基準であり、$3'$ への赤い線で寄与する。ラベル付けされたノードは集合 $S,S'$ にグループ化でき、$\fD,\gD$ を修正して異なるグラフを得ることにより、より高いフローを得るためにより良く利用できるノードを特定する。証明はプロット (d) に記述された部分的なフローの和を含む。}
\end{figure}

双対上昇法は、実行可能性を保証し目的関数を改善する集合 $S,S'$ によって生成される任意のベクトルで実行可能な解 $(\fD,\gD)$ を修正することにより進む。集合 $S,S'$ が命題~\ref{prop-dual-ascent-dir}の証明で構成的に与えられたものであり、ステップサイズ $\varepsilon$ が命題~\ref{prop-feasibility-dual-pair}の証明で定義されたものである場合、\emph{主双対}法として知られる方法が復元される。この方法はマッチング問題に対するハンガリアン法に帰着する。双対上昇法はネットワークシンプレックスの双対版と類似点を共有するが、少なくとも2つの重要な点で異なる。シンプレックス型の方法は、現在の解が実行可能集合（双対に対しては $\PotentialsD(\C)$）の\emph{端点}であることを常に保証するが、ここで提示された双対上昇法はそのような仮定を置かず、実行可能集合の内部にある反復点を自由に生成できる。さらに、双対ネットワークシンプレックスが $(\fD,\gD)$ を修正して線形（周辺制約）を満たすが非負性は収束時にのみ満たす主解 $\P$ を生成するのに対し、双対上昇法は常に非負だが必ずしも周辺制約を満たさない主解 $\P$ を代わりに構築する。

\section{オークションアルゴリズム}\label{s-auction}
オークションアルゴリズムは~\citet{bertsekas1981new}によって最初に提案され、後に~\citep{bertsekas1988dual}で改良された。このアルゴリズムのいくつかの経済学的解釈が提案されている（例えば \citet{bertsekas1992auction}を参照）。このアルゴリズムは任意の周辺に対して適用できるが、ここでは最適割当問題を解くための定式化で提示する。

\paragraph{相補スラック性。} 最適割当問題では、optimal transport 問題に対して提示された主双対条件がより簡単に定式化できることに注意せよ。任意の端点的解 $\P$ は必然的にある $\si$ に対する置換行列 $\P_\si$ である（式~\eqref{prop-primal-dual-optimality}を参照）。主変数 $\P_{\si^\star}$ と双対最適解 $\fD^\star,\gD^\star$ が与えられたとき、必然的に
$$
\fD^\star_i + \gD^\star_{\si^\star_i} = \C_{i,{\si_i^\star}}
$$
が成り立つ。また、第~\ref{sec-c-transforms}節で述べた $\C$-変換の原理により、$\fD^\star$ を $\gD^{\bar\C}$ に等しく選べることも思い出そう。したがって
\eql{\label{eq-dual-conditions-assign-1}
\C_{i,\si^\star_i} - \gD^\star_{\si_i}  = \min_{j} \C_{i,j} - \gD^\star_{j}
}
が成り立つ。逆に、あるベクトル $\gD$ と置換 $\sigma$ が存在して
\eql{\label{eq-dual-conditions-assign}
\C_{i,\si_i} - \gD_{\si_i}  = \min_{j} \C_{i,j} - \gD_{j}
}
が成り立つならば、$\si$ が最適割当であり $\gD^{\bar\C},\gD$ が最適双対ペアであるという意味で、両者は最適であることを示すのは容易である。


\paragraph{部分割当と $\varepsilon$-相補スラック性。} オークションアルゴリズムの目標は、三つ組 $S,\xi,\gD$ を反復的に修正することである。ここで $S$ は $\range{n}$ の部分集合、$\xi$ は部分割当ベクトル、すなわち $S$ から $\range{n}$ への単射、$\gD$ は双対ベクトルである。双対ベクトルは\emph{近似的}相補スラック性~\eqref{eq-dual-conditions-assign}を満たす解に収束することが意図されており、$S$ は $\xi$ が置換を記述するに従って $\range{n}$ を覆うように成長する。アルゴリズムは各反復後に以下の3つの性質を維持することにより動作する：
\begin{enumerate}[label={(\alph*)}]
	\item $\forall i\in S,\quad \C_{i,\xi_i} - \gD_{\xi_i}  \leq \varepsilon + \min_{j} \C_{i,j} - \gD_{j}$（$\varepsilon$-CS）。
	\item $S$ のサイズは各反復で増加するのみである。
	\item ある添字 $i$ が存在して $\gD_i$ が少なくとも $\varepsilon$ だけ減少する。
\end{enumerate}

\paragraph{オークションアルゴリズムの更新。} 点 $j$ が与えられたとき、オークションアルゴリズムは通常の $\C$-変換に現れる最適値だけでなく、2番目に良い値も使用する：
$$j^1_i \in \argmin_{j} \C_{i,j} - \gD_{j}, \quad j^2_i \in \argmin_{j\ne j^1_i} \C_{i,j} - \gD_{j}, $$
これを用いて、$i\notin S$ である添字 $i$ に対する $\gD$ の以下の更新、および $S$ と $\xi$ の更新を定義する：
\begin{enumerate}\item \textbf{$\gD$ の更新}：$\gD$ の $j^1_i$ 番目の成分から、$\varepsilon$ と調整コスト $\{\C_{i,j}-\gD_j\}_j$ の最低値と2番目に低い値の差の和を引く。
\eql{
\begin{aligned}
	\gD_{j^1_{i}} &\leftarrow \gD_{j^1_{i}} - \underbrace{\left((\C_{i,j^2_{i}}-\gD_{j^2_{i}})-(\C_{i,j^1_{i}}-\gD_{j^1_{i}})+\varepsilon\right)}_{\geq \varepsilon >0}\label{eq-auction-update}\\
	&= \C_{i,j^1_{i}} - (\C_{i,j^2_{i}}-\gD_{j^2_{i}}) -\varepsilon.\\
\end{aligned}
}
\item \textbf{$S$ と $\xi$ の更新}：$\xi_{i'}=j^1_{i}$ となる添字 $i'\in S$ が存在すれば、$S\leftarrow S\setminus \{i'\}$ として除去する。$\xi_{i}=j^1_{i}$ と設定し、$i$ を $S$ に追加する。$S\leftarrow S \cup \{i\}$。
\end{enumerate}

\paragraph{アルゴリズム的性質。} アルゴリズムは割当済み点の空集合 $S=\emptyset$（割当なし、部分割当ベクトル $\xi$ は空）と $\gD=\zeros_n$ から開始し、$S=\range{n}$ となったら終了し、終了するまで上記の両ステップを繰り返す。性質 {(b)} と {(c)} が各反復後に有効であることは、更新の性質から明らかである（式~\eqref{eq-auction-update}を見れば十分である）。$\varepsilon$-相補スラック性は、$S=\emptyset$ であるため最初の反復では容易に満たされる。反復がこの性質を保存することは次の命題で示される。
\begin{prop} オークションアルゴリズムは各反復で $\varepsilon$-相補スラック性を維持する。
\end{prop}
\begin{proof}
	$\gD,\xi,S$ をある反復の開始時の3つの変数とする。したがって、任意の $i'\in S$ に対して
	$$\C_{i,\xi_{i'}} - \gD_{\xi_{i'}}  \leq \varepsilon + \min_{j} \C_{{i'},j} - \gD_{j} $$
が成り立つと仮定する。ある反復で考えられる特定の $i\notin S$ について考える。3つの更新が行われる：$\gD,\xi,S$ は添字 $j^1_i$ と $j^2_i$ を用いて $\gD^{\text{n}},\xi^{\text{n}},S^{\text{n}}$ に更新される。より正確には、$\gD^{\text{n}}$ は要素 $j^1_i$ を除いて $\gD$ に等しく、その値は
$$\gD_{j^1_{i}}^\text{n} = \gD_{j^1_{i}} - \left((\C_{i,j^2_{i}}-\gD_{j^2_{i}})-(\C_{i,j^1_{i}}-\gD_{j^1_{i}})\right)-\varepsilon \leq \gD_{j^1_{i}}-\varepsilon$$
であり、$\xi^{\text{n}}$ は $i$ 番目の要素が $j^1_i$ に等しいことを除いて $\xi$ に等しく、$S^{\text{n}}$ は $\{i\}$ と $S$ の和集合に等しい（1つの要素が除去される可能性がある）。$\gD^{\text{n}}$ の更新は次のように書き直せる：
$$\gD^{\text{n}}_{j^1_i}= \C_{i,j^1_{i}} - (\C_{i,j^2_{i}}-\gD_{j^2_{i}}) -\varepsilon;$$
したがって
$$\C_{i,j^1_{i}} - \gD^{\text{n}}_{j^1_i}=  \varepsilon + (\C_{i,j^2_{i}}-\gD_{j^2_{i}}) = \varepsilon+\min_{j\ne j^1_i} (\C_{i,j}-\gD_{j}).$$
が成り立つ。$-\gD\leq -\gD^{\text{n}}$ であるから、これは
$$\C_{i,j^1_{i}} - \gD^{\text{n}}_{j^1_i} = \varepsilon + \min_{j\ne j^1_i} (\C_{i,j}-\gD_{j})\leq \varepsilon+\min_{j\ne j^1_i} (\C_{i,j}-\gD^{\text{n}}_{j})$$
を含意し、不等式は $j=j^1_i$ に対しても明らかに成り立つため、添字 $i$ に対する $\varepsilon$-相補スラック性が得られる。他の添字 $i'\ne i$ については、$\gD^{\text{n}}\leq\gD$ であるから、
$$\C_{i,\xi^{\text{n}}_{i'}} - \gD^{\text{n}}_{\xi^{\text{n}}_{i'}} = \C_{i,\xi_{i'}} - \gD_{\xi_{i'}} \leq \varepsilon + \min_{j} \C_{{i'},j} - \gD_{j} \leq \varepsilon + \min_{j} \C_{{i'},j} - \gD^{n}_{j} $$
という不等式の列が成り立つ。
\end{proof}

\begin{prop} オークションアルゴリズムのステップ数は最大 $N=n\|\C\|_\infty/\varepsilon$ である。
\end{prop}
\begin{proof}アルゴリズムが $T>N$ ステップ後に停止していないと仮定する。するとある添字 $j$ が存在して、$\xi$ の像に含まれない、すなわち価格座標 $\gD_j$ が一度も更新されておらず依然として $\gD_j=0$ である。この場合、$\gD_{j'}$ が $n>\|\C\|_\infty/\varepsilon$ 回更新されるような添字 $j'$ は存在し得ない。実際、もしそうであれば、任意の添字 $i$ に対して
$$\gD_{j'}\leq-n\varepsilon<-\|\C\|_\infty\leq -\C_{i,j} = \gD_j-\C_{i,j},$$
が成り立ち、すべての $i$ に対して
$$\C_{i,j'}-\gD_{j'}> \C_{i,j}+ (\C_{i,j}-\gD_j)$$
となり、$\varepsilon$-CS に矛盾する。したがって、各変数に対する更新は $\|C\|_\infty/\varepsilon$ 回を超えることはできず、反復の総数 $T$ は $n\|\C\|_\infty/\varepsilon=N$ より大きくなり得ない。
\end{proof}

\begin{rem} この結果はアルゴリズムが終了するための素朴な操作数 $N^3\|\C\|_\infty/\varepsilon$ を与える。この計算量は $\varepsilon$-スケーリングとして知られる巧妙な方法を用いることで $N^3 \log\|\C\|_\infty$ に削減できる。これは各反復で $\varepsilon$ の値を減少させるように設計されている（\cite[p. 264]{bertsekas1998network}）。
\end{rem}

\begin{prop} オークションアルゴリズムはコストが $n\varepsilon$ だけ準最適な割当を見つける。
\end{prop}
\begin{proof}
	$\sigma,\gD^\star$ を割当問題のコスト行列 $\C$ の主最適解と双対最適解とし、最適値を $$t^\star=\sum \C_{i,\sigma_i} = \sum_i \min_{j}\C_{i,j}-\gD^\star_j+\sum_j\gD^\star_j$$ とする。
$\xi,\gD$ を終了時にオークションアルゴリズムが出力する解とする。$\varepsilon$-CS 条件より、任意の $i\in S$ に対して
$$\min_{j} \C_{i,j} - \gD_{j}\geq \C_{i,\xi_i} - \gD_{\xi_i} -\varepsilon$$
が成り立つ。したがって $\gD$ の単純な準最適性により、
	$$\begin{aligned}t^\star&\geq \sum_i \left(\min_{j} \C_{i,j}-\gD_{j}\right) + \sum_j \gD_j\\
&\geq
	\sum_i -\varepsilon+\left(\C_{i,\xi_i}-\gD_{\xi_i}\right) + \sum_j \gD_j = -n\varepsilon + \sum_i \C_{i,\xi_j} \geq  - n\varepsilon + t^\star.\end{aligned},$$
が得られる。ここで第二の不等式は $\varepsilon$-CS から、次の等式は $\gD_{\xi_i}$ の項と $\gD_{j}$ の項の打ち消しから、最後の不等式は置換としての $\xi$ の準最適性から従う。
\end{proof}


オークションアルゴリズムは、$\C$-変換の機構を使用する代替的な方法と見なすことができる。次に、正則化に基づく別のアプローチ、いわゆる Sinkhorn アルゴリズムを探求する。これは~\citep{schmitzer2016stabilized}で議論されているように、オークションアルゴリズムとも類似点を持つ。

最後に、ユークリッド空間の低次元正則グリッド上では、これらの古典的な線形ソルバーをマルチスケール戦略と結合して、顕著な高速化を得ることが可能であることを注記する~\citep{schmitzer2016sparse,oberman2015efficient}。
